% TOP_PROBLEM: {"id": 3100009, "title": "What is the minimum number of n-simplexes needed to triangulate the n-cube", "category_id": 2, "category": {"id": 2, "name": "combinatorics", "display_name": "Combinatorics", "description": "Counting problems, graph theory, discrete structures.", "slug": "combinatorics", "order_index": 2, "created_at": "2026-07-31T15:26:25.671Z"}, "status": "open", "problem_number": "AMR-030-0009", "set_id": 13, "statusQualification": "Uses the standard vertex-triangulation convention for cube simplexity; Steiner triangulations and non-face-to-face dissections are separate variants."}
% FIRST_POSTED: 2026-10-02
% ENTRY_KIND: statement_only
% UnsolvedMath ID 3100009; source catalogue code AMR-030-0009
% !TeX program = lualatex
% Definition and sources only; this file is not an LLM solution attempt.
\documentclass[11pt,a4paper]{article}
\usepackage[margin=25mm]{geometry}
\usepackage{fontspec}
\setmainfont{lmroman10-regular.otf}[BoldFont=lmroman10-bold.otf,
 ItalicFont=lmroman10-italic.otf,BoldItalicFont=lmroman10-bolditalic.otf]
\usepackage{amsmath,amssymb,mathtools}
\usepackage{unicode-math}
\setmathfont{latinmodern-math.otf}
\AtBeginDocument{\let\setminus\smallsetminus}
\usepackage{xurl,hyperref}
\hypersetup{colorlinks=true,urlcolor=blue}
\setcounter{secnumdepth}{0}
\setlength{\parindent}{0pt}
\setlength{\parskip}{5pt}
\setlength{\emergencystretch}{3em}
\begin{document}
\section*{What is the minimum number of n-simplexes needed to triangulate the n-cube}\label{3100009}
{\small UnsolvedMath ID 3100009 · AMR-030-0009 · Combinatorics · AMR Open Problem Lists}
\subsection{Definitions and mathematical statement}
Let $Q_n=[0,1]^n$ and let its vertex set be $\{0,1\}^n$.
A vertex triangulation of $Q_n$ is a finite collection of
nondegenerate $n$-simplices whose vertices belong to $\{0,1\}^n$,
whose union is $Q_n$, and such that the intersection of any two
simplices is either empty or a common face of both. In particular,
their interiors are disjoint. Lower-dimensional faces are included
implicitly in the resulting simplicial complex.

Define the simplexity in this convention by
\[
 t(n)=\min\{\text{number of full-dimensional simplices in such a
 triangulation}\}.
\]
Determine $t(n)$, or obtain asymptotically sharp bounds as
$n\to\infty$. The usual triangulation obtained by ordering the
coordinates shows $t(n)\le n!$, but it need not be optimal.

The corner restriction excludes added Steiner vertices, and the
common-face condition excludes arbitrary simplicial dissections
with incompatible subdivisions across interfaces. These are
different minimization problems, even when useful bounds compare
them. Counting refers only to $n$-dimensional simplices, not to all
faces in the complex. An upper bound requires an actual compatible
triangulation covering the cube, while a lower bound must apply to
every such triangulation, irrespective of symmetry or regularity.
\subsection{Short English statement}
Determine the minimum number of n-simplices in a triangulation of the n-cube using only cube vertices.
\subsection{Sources}
\begin{itemize}
\item[S1] ULAM.AI and the UnsolvedMath Contributors, supplied problems.json, record 3100009 (AMR-030-0009), AMR Open Problem Lists. Catalogue identity and source transcription; CC BY 4.0.
\url{https://huggingface.co/datasets/ulamai/UnsolvedMath}.
\item[S2] Cooper - Combinatorial Problems I Like (2020 snapshot), Geometry, source-order bullet 9. Original source indexed by AMR-030-0009.
\url{https://people.math.sc.edu/cooper/combprob.html}.
\item[S3] Hughes, Simplexity of the cube, definitions of triangulation and simplexity.
\url{https://www.math.ucdavis.edu/~deloera/MISC/LA-BIBLIO/trunk/Hughes/Hughes6.pdf}.
\end{itemize}
\end{document}
